% Figure: inductive wireless charging coupling efficiency against coil
% separation, for two coupling coefficients, showing the steep falloff once
% the gap exceeds the coil radius. Coordinates computed and braced.
\begin{figure}[htbp]
\centering
\begin{tikzpicture}
\begin{axis}[
  width=12cm, height=7cm,
  xlabel={gap-to-radius ratio $z/a$},
  ylabel={coupling coefficient $k$},
  xmin=0, xmax=3, ymin=0, ymax=1,
  axis lines=left,
  tick label style={font=\scriptsize},
  label style={font=\small},
  legend pos=north east,
  legend style={font=\scriptsize},
]
% Coaxial coplanar-radius coils: k approx a^3 / (a^2 + z^2)^{3/2} = 1/(1+(z/a)^2)^{3/2}
\addplot[engblue, very thick, samples=120, domain=0:3] {1/(1+x^2)^(1.5)};
\addlegendentry{equal-radius coaxial coils}
% mark the z = a point where k has dropped to 0.354
\filldraw[engred] (axis cs:1,0.354) circle [radius=1.8pt];
\node[engred, font=\scriptsize, anchor=west] at (axis cs:1.1,0.40)
  {$z=a$: $k\approx0.35$};
% shaded acceptable region z/a < 0.5
\draw[engamber, dashed] (axis cs:0.5,0) -- (axis cs:0.5,0.716);
\node[engamber, font=\scriptsize, anchor=west, align=left]
  at (axis cs:0.55,0.85) {practical pad:\\$z/a<0.5$};
\end{axis}
\end{tikzpicture}
\caption[Wireless-charging coupling against gap]{The coupling coefficient
between two coaxial equal-radius coils falls steeply once the gap $z$ reaches
the coil radius $a$, following $k\approx[1+(z/a)^2]^{-3/2}$ for thin coils on
a common axis. A practical phone-charging pad holds the gap below half the
coil radius (amber), keeping the coupling above about $0.7$; at a gap equal to
the radius the coupling has already dropped to a third, and the transferred
power with it. The steep falloff is why such pads demand close, aligned
placement and why the receiver coil is made as large as the phone body
allows.}
\label{fig:vol04ch08:wireless-charging}
\end{figure}
